% Transcription of Cayley's Collected Mathematical Papers, Vol. 10, pages 451--500
% Papers 702, 703, 704

\documentclass[12pt,a4paper]{article}

% ==================== PACKAGES ====================
\usepackage[utf8]{inputenc}
\usepackage[T1]{fontenc}
\usepackage[english]{babel}
\usepackage{amsmath,amssymb,amsthm}
\usepackage{mathtools}
\usepackage{bm}
\usepackage{eufrak}  % Fraktur/mathfrak fonts
\usepackage{array}
\usepackage{booktabs}
\usepackage{longtable}
\usepackage{geometry}
\usepackage{fancyhdr}
\usepackage{titlesec}
\usepackage{enumitem}
\usepackage{hyperref}

% ==================== GEOMETRY ====================
\geometry{
    paperwidth=6.125in,
    paperheight=9.25in,
    textwidth=4.5in,
    textheight=7.25in,
    top=0.875in,
    inner=0.8125in,
    outer=0.8125in,
    bottom=0.875in,
    headheight=12pt,
    footskip=24pt
}

% ==================== PAGE STYLE ====================
\pagestyle{fancy}
\fancyhf{}
\fancyhead[LE]{\thepage}
\fancyhead[RO]{\thepage}
\renewcommand{\headrulewidth}{0pt}
\renewcommand{\footrulewidth}{0pt}

% ==================== CUSTOM COMMANDS ====================
\DeclareMathOperator{\expfc}{exp.}

% Underlined numbers for characteristics
\newcommand{\ul}[1]{\underline{#1}}

% Elliptic function shortcuts
\DeclareMathOperator{\sn}{sn}
\DeclareMathOperator{\cn}{cn}
\DeclareMathOperator{\dn}{dn}
\DeclareMathOperator{\Imop}{\Im}  % alternation operator (using Im symbol)
\newcommand{\cbtw}{\vert}  % Cayley's between notation (renamed to avoid conflict with amssymb)

% ==================== SECTION FORMATTING ====================
\titleformat{\section}{\normalfont\large\bfseries\centering}{\thesection.}{0.5em}{}
\titleformat{\subsection}{\normalfont\normalsize\bfseries\centering}{\thesubsection.}{0.5em}{}
\titlespacing*{\section}{0pt}{1.5ex plus 0.5ex minus 0.2ex}{1ex plus 0.2ex}
\titlespacing*{\subsection}{0pt}{1ex plus 0.3ex minus 0.2ex}{0.5ex plus 0.2ex}

% ==================== THEOREM STYLES ====================
\theoremstyle{definition}
\newtheorem*{remark}{Remark}

% ==================== DOCUMENT ====================
\begin{document}
\setlength{\parindent}{1.5em}
\setlength{\parskip}{0pt plus 1pt}

% ============================================================
% PAPER 702: ON THE TRIPLE theta-FUNCTIONS (continued)
% ============================================================
\thispagestyle{empty}
\begin{center}
\large\bfseries ON THE TRIPLE $\vartheta$-FUNCTIONS.\par\vspace{0.5em}
\normalsize\normalfont [702
\end{center}

\vspace{1em}

\noindent\hangindent=2em even characteristic $(56 \times 5 =)\ 280$ tripairs, which are however $70$ tripairs each taken $4$~times. A tripair contains in all $(2^3 =)\ 8$ hemi-tripairs, but these divide themselves into two sets each of $4$ hemi-tripairs such that for each hemi-tripair of the first set the three characteristics have a given sum, and for each hemi-tripair of the second set the three characteristics have a different given sum. Hence considering the $70$ tripairs corresponding as above to a given even characteristic, in any one of the $70$ tripairs, there is a set of $4$ hemi-tripairs such that in each of them the sum of the three characteristics is equal to the given even characteristic; and taking the bitangents $f$, $g$, $h$ to correspond to any one of these hemi-tripairs, the bitangents which correspond to the other three hemi-tripairs will be $b$, $c$, $f$; $c$, $a$, $g$ and $a$, $b$, $h$ respectively; and we thus obtain from any one of these one and the same representation
\[
\alpha\sqrt{bcf} + \beta\sqrt{cag} + \gamma\sqrt{abh} + \delta\sqrt{fgh} = 0
\]
of the cubic hexad. And the $70$ tripairs give thus the $70$ representations of the same cubic hexad.

The whole number of hemi-tripairs is $36 \times 56 = 2016$: it may be remarked that there exists a system of $288$ heptads, each of $7$ odd characteristics such that selecting at pleasure any $3$ characteristics out of the heptad, we obtain always a hemi-tripair: we have thus in all $288 \times 35 = 10080$ hemi-tripairs: this is $= 2016 \times 5$, or we have the $2016$ hemi-tripairs each taken $5$ times. Weber's Table~II.\ exhibits $36$ out of the $288$ heptads.

I recall that in the algorithm derived from Hesse's theory the bitangents are represented by the duads $12$, $13$, \dots, $78$ formed with the eight figures $1$, $2$, $3$, $4$, $5$, $6$, $7$, $8$; these duads correspond to the odd characteristics as shown in the Table, and the table shows also triads corresponding to all the even characteristics except $\begin{matrix} 000 \\ 000 \end{matrix}$.

\begin{center}
\small
\textit{Top line of characteristic.}
\end{center}

\begin{center}
\small
\begin{tabular}{|c|c|c|c|c|c|c|c|c|}
\hline
 & 000 & 100 & 010 & 110 & 001 & 101 & 011 & 111 \\
\hline
000 &  & 236 & 345 & 137 & 467 & 156 & 124 & 257 \\
\hline
100 & 237 & 67 & 136 & 12 & 157 & 48 & 256 & 35 \\
\hline
010 & 245 & 127 & 23 & 68 & 134 & 357 & 15 & 47 \\
\hline
110 & 126 & 13 & 78 & 145 & 356 & 25 & 46 & 234 \\
\hline
001 & 567 & 146 & 125 & 247 & 45 & 17 & 38 & 26 \\
\hline
101 & 147 & 58 & 246 & 34 & 16 & 123 & 27 & 367 \\
\hline
011 & 135 & 347 & 14 & 57 & 17 & 36 & 167 & 456 \\
\hline
111 & 346 & 24 & 56 & 235 & 37 & 267 & 457 & 18 \\
\hline
\end{tabular}
\end{center}

\begin{center}
\small
\textit{Bottom line of characteristic.}
\end{center}

\vspace{1em}

See my ``Algorithm of the triple $\vartheta$-functions,'' Crelle, t.~LXXXVII.\ p.~165, [701].

\vspace{1em}

\noindent\textbf{451.}\quad
The $(35 + 28 =)\ 63$ hexpairs then are

$35$ hexpairs such as

\noindent\hangindent=3em say this is \textit{1234.5678} or for shortness $567$ (the $8$ going always with the expressed triad):

\noindent that is, $567$ denotes the hexpair
\[
12.34;\ 13.24;\ 14.23;\ 56.78;\ 57.68;\ 58.67:
\]

\noindent and $28$ hexpairs such as

\noindent\hangindent=3em say this is $12$; that is, $12$ denotes the hexpair
\[
13.32;\ 14.42;\ 15.52;\ 16.62;\ 17.72;\ 18.82.
\]

It is to be noticed that the odd characteristics, as represented by their duad symbols, can be added by the formulae

$12 + 23 = 13$, etc.,

or, what is the same thing,

$12 + 13 + 23 = 0$, etc.,

and

$12 + 34 = 13 + 24 = 14 + 23 = 56 + 78 = 57 + 68 = 58 + 67 = 567$, etc.

Thus, referring to the table,

\noindent\hangindent=3em $12 + 23 = 13$ means
$\begin{matrix} 100 \\ 010 \end{matrix} +
\begin{matrix} 010 \\ 100 \end{matrix} =
\begin{matrix} 110 \\ 110 \end{matrix}$;

\noindent and

\noindent\hangindent=3em $12 + 34 = 567$ means
$\begin{matrix} 100 \\ 010 \end{matrix} +
\begin{matrix} 101 \\ 100 \end{matrix} =
\begin{matrix} 001 \\ 110 \end{matrix}$,

\noindent which are right.

The $288$ heptads are

\noindent $8$ heptads such as
\[
\begin{matrix} 1, & 2, & 3, & 4, & 5, & 6, & 7, & 8 \\ & 12, & 13, & 14, & 15, & 16, & 17, & 18 \end{matrix}
\]
say this is the heptad $1$, denoting the seven duads $12$, $13$, $14$, $15$, $16$, $17$, $18$;

and $280$ heptads such as
\[
\begin{matrix} 1, & 2, & 3, & 4, & / & 5, & \backslash & 6, & 7, & 8 \\ & 12, & 13, & 14, & 15, & & 67, & 68, & 78 \end{matrix}
\]
say this is the heptad $1.678$, denoting the seven duads $12$, $13$, $14$, $15$, $67$, $68$, $78$.

We hence see that the $2016$ hemi-tripairs are

\noindent (I.)\quad $280$ hemi-tripairs
\[
\begin{matrix} 12, & 13, & 23 \end{matrix},
\]
say this is $12$ (hemi-tripair), denoting the three duads $12$, $13$, $23$;

\noindent (II.)\quad $1680$ hemi-tripairs
\[
\begin{matrix} 12, & 67, & 68 \end{matrix},
\]
say this is $12.678$, denoting the three duads $12$, $67$, $68$;

\noindent (III.)\quad $56$ hemi-tripairs
\[
\begin{matrix} 123, & 234, & 134 \end{matrix},
\]
say this is $123$, denoting the three duads $123$, $234$, $134$;

\noindent in all $280 + 1680 + 56 = 2016$.

We further see how each hemi-tripair may be completed into a tripair in different ways: thus (I.) gives the $5$ tripairs
\[
\begin{array}{c|c|c|c|c}
12 & 13 & 14 & 15 & 16 \\
\hline
34 & 24 & 23 & 23 & 23 \\
56 & 56 & 56 & 46 & 45 \\
78 & 78 & 78 & 78 & 78
\end{array}
\]
the last duad being $5$, $6$, $7$ or $8$; (II.) gives the $3$ tripairs
\[
\begin{array}{c|c|c}
12 & 12 & 12 \\
\hline
67 & 68 & 78 \\
34 & 34 & 34 \\
(/ & 7)8 & 6\\
4 & 4 & 4
\end{array}
\]
and the $2$ tripairs
\[
\begin{array}{c|c}
12 & 12 \\
\hline
67 & 68 \\
(/ & 5\\
5 & 5
\end{array}
\]
while (III.) gives the $5$ tripairs
\[
\begin{array}{c|c|c|c|c}
123 & 123 & 123 & 123 & 123 \\
\hline
456 & 456 & 457 & 457 & 467 \\
78 & 78 & 68 & 6 & 5\\
 & & & 8 & 8
\end{array}
\]

\vspace{1em}

\begin{center}
\textbf{452}
\end{center}

\vspace{0.5em}

\noindent\textbf{[702}

\vspace{0.5em}

To each even characteristic there belongs a system of $56$ tripairs: thus for $\begin{matrix} 000 \\ 000 \end{matrix}$, the characteristic $\begin{matrix} 000 \\ 000 \end{matrix}$, the $56$ hemi-tripairs are $123$, that is, $12$, $13$, $23$, etc.; and in the tripair, $1234$, that is, $12.34$; $13.24$; $14.23$, we have the set of four hemi-tripairs each of which the sum of the three characteristics is $= \begin{matrix} 000 \\ 000 \end{matrix}$:
\[
\begin{matrix} 000 \\ 000 \end{matrix}
\begin{cases}
(12 + 23 + 13 = \begin{matrix} 000 \\ 000 \end{matrix}, \text{ etc.}):
\end{cases}
\]
whence the other $70$ hemi-tripairs: for any such even characteristic, the tripairs are

$2$ such as $1234$,

$16$ such as $54(123)$,

$16$ such as $15(678)$,

$36$ such as $(5162)34.78$;

\noindent in all $2 + 16 + 16 + 36 = 70$; and in each of these it is easy to select the hemi-tripairs for which the sum of the $3$ duads is $= 567$.

\vspace{1em}

\begin{flushright}
Cambridge, 27 December, 1878.
\end{flushright}

\vspace{2em}

% ============================================================
% PAPER 703: ON THE ADDITION OF THE DOUBLE theta-FUNCTIONS
% ============================================================
\begin{center}
\large\bfseries 703.\par\vspace{0.3em}
ON THE ADDITION OF THE DOUBLE $\vartheta$-FUNCTIONS.\par\vspace{0.5em}
\normalsize\normalfont
[From the Journal f\"{u}r die reine und angewandte Mathematik (Crelle), t.~LXXXVIII.\ (1879), pp.~74---81.]
\end{center}

\vspace{1em}

\noindent\textbf{455.}\quad
I assume in general
\[
\xi = a - \theta . b - \theta . c - \theta . d - \theta . e - \theta . f - \theta,
\]
and I consider the variables $x$, $y$, $z$, $w$, $p$, $q$, connected by the equations
\begin{align*}
\frac{1}{\sqrt{X}},\ \frac{1}{\sqrt{Y}},\ \frac{1}{\sqrt{Z}},\ \frac{1}{\sqrt{W}},\ \frac{1}{\sqrt{P}},\ \frac{1}{\sqrt{Q}} &=
\frac{x}{\sqrt{X}},\ \frac{y}{\sqrt{Y}},\ \frac{z}{\sqrt{Z}},\ \frac{w}{\sqrt{W}},\ \frac{p}{\sqrt{P}},\ \frac{q}{\sqrt{Q}} \\
&= \frac{x^2}{\sqrt{X}},\ \frac{y^2}{\sqrt{Y}},\ \frac{z^2}{\sqrt{Z}},\ \frac{w^2}{\sqrt{W}},\ \frac{p^2}{\sqrt{P}},\ \frac{q^2}{\sqrt{Q}} \\
&= \frac{\sqrt{X'}}{\sqrt{X}},\ \frac{\sqrt{Y'}}{\sqrt{Y}},\ \frac{\sqrt{Z'}}{\sqrt{Z}},\ \frac{\sqrt{W'}}{\sqrt{W}},\ \frac{\sqrt{P'}}{\sqrt{P}},\ \frac{\sqrt{Q'}}{\sqrt{Q}},
\end{align*}
equivalent to two independent equations, which serve to determine $p$, $q$, or say the symmetrical functions $p + q$ and $pq$, in terms of $x$, $y$, $z$, $w$.

These equations, it is well known, constitute a particular integral of the differential equations
\[
\frac{dx}{\sqrt{X}} = \frac{dy}{\sqrt{Y}} = \frac{dz}{\sqrt{Z}} = \frac{dw}{\sqrt{W}} = \frac{dp}{\sqrt{P}} = \frac{dq}{\sqrt{Q}},
\]
or what is the same thing, regarding $\xi$, $\xi'$ as arbitrary constants, they constitute the general integral of the differential equations
\[
\frac{dx}{\sqrt{X}} = \frac{dy}{\sqrt{Y}} = \frac{dz}{\sqrt{Z}} = \frac{dw}{\sqrt{W}} = \frac{x\,dx}{\sqrt{X}} = \frac{y\,dy}{\sqrt{Y}} = \frac{z\,dz}{\sqrt{Z}} = \frac{w\,dw}{\sqrt{W}}.
\]

\vspace{1em}

\begin{center}
\textbf{456}
\end{center}

\vspace{0.5em}

\noindent\textbf{[703}

\vspace{0.5em}

I attach the numbers $1$, $4$, $3$, $2$, $5$, $6$ to the variables $x$, $y$, $z$, $w$, $p$, $q$, respectively: and write

\noindent\hangindent=3em (six equations),
\[
A_{12} = \sqrt{a - x . a - y},\quad A_{34} = \sqrt{a - z . a - w},\quad A_{56} = \sqrt{a - p . a - q};
\]

\noindent\hangindent=3em (ten equations),
\[
AB_{12} = \sqrt{a - x . b - x . c - x . d - x . e - x . f - x . a - y . b - y . c - y . d - y . e - y . f - y}; \text{ etc.}
\]

where it is to be borne in mind that $AB$ is an abbreviation for $ABF.CDE$, and so in all other cases, the letter $F$ belonging always to the expressed duad: there are sixteen functions $A_{12}$, $A_{34}$, $A_{56}$, $AB_{12}$, $AC_{12}$, $AD_{12}$, $AE_{12}$, $AF_{12}$, $BC_{12}$, $BD_{12}$, $BE_{12}$, $BF_{12}$, $CD_{12}$, $CE_{12}$, $CF_{12}$, $DE_{12}$ being functions of $x$ and $y$, of $z$ and $w$, and of $p$ and $q$, according as the suffix is $12$, $34$, or $56$.

It is to be shown that the $16$ functions $A_{12}$, $AB_{12}$, etc., of $x$ and $y$, of the addition of the double $\vartheta$-functions. I use when convenient the abbreviated notations
\[
a - x = a_1,\quad a - y = a_2, \text{ etc.},
\]
\[
b_{12} = x - y,\quad c_{31} = z - x, \text{ etc.},
\]
\[
e_{31} = z - w,\quad f_{12} = p - q;
\]
we have of course
\[
X = a_1 b_1 c_1 d_1 e_1 f_1,
\]
\[
A_{12} = \sqrt{a_1 a_2 b_1 b_2 c_1 c_2 d_1 d_2 e_1 e_2 f_1 f_2} = \sqrt{a_1 a_2 b_1 b_2 c_1 c_2 d_1 d_2 e_1 e_2} \cdot f_{12}, \text{ etc.}
\]

\vspace{1em}

\noindent\textbf{17.}\quad
Proceeding to the investigation, the six equations obtained by writing therein $\alpha$, $\beta$, $\gamma$, $\delta$, $\epsilon$ for the values $x$, $y$, $z$, $w$, $p$, $q$ successively; the four equations
\begin{align*}
\alpha x^2 + \beta x + \gamma + \delta/x &= \epsilon\sqrt{X}, \\
\alpha y^2 + \beta y + \gamma + \delta/y &= \epsilon\sqrt{Y}, \\
\alpha z^2 + \beta z + \gamma + \delta/z &= \epsilon\sqrt{Z}, \\
\alpha w^2 + \beta w + \gamma + \delta/w &= \epsilon\sqrt{W},
\end{align*}
serving to determine the ratios $\alpha : \beta : \gamma : \delta : \epsilon$ in terms of $x$, $y$, $z$, $w$; and we have then for the determination of $p$, $q$ the remaining two equations
\begin{align*}
\alpha p^2 + \beta p + \gamma + \delta/p &= \epsilon\sqrt{P}, \\
\alpha q^2 + \beta q + \gamma + \delta/q &= \epsilon\sqrt{Q},
\end{align*}
which two equations may be replaced by the identity
\[
(\alpha\theta^2 + \beta\theta + \gamma + \delta/\theta)^2 - \epsilon^2\xi = a^2 - \epsilon^2 \cdot \theta - x . \theta - y . \theta - z . \theta - w . \theta - p . \theta - q.
\]

\vspace{1em}

\begin{center}
\textbf{457}
\end{center}

\vspace{0.5em}

\noindent\textbf{[703}

\vspace{0.5em}

Writing herein $\theta = $ any one of the values $a$, $b$, $c$, $d$, $e$, $f$, for instance $\theta = a$, and taking the square root of each side, we have
\[
\alpha a^2 + \beta a + \gamma + \delta/a = \sqrt{a^2 - \xi}\sqrt{a - x . a - y . a - z . a - w . a - p . a - q},
\]
or as this may be written
\[
\alpha a^2 + \beta a^2 + \gamma a + \delta = \sqrt{a^2 - \xi} \cdot A_{12} A_{34} A_{56},
\]
which equation when reduced gives the proportional value of $A_{12}$.

For the reduction we require the value of $\alpha a^2 + \beta a^2 + \gamma a + \delta$: calling this for the moment $\Pi$, we join to the four equations a fifth equation
\[
\alpha a^2 + \beta a^2 + \gamma a + \delta = \Pi.
\]

Eliminating $\alpha$, $\beta$, $\gamma$, $\delta$, we find
\[
\begin{vmatrix}
\sqrt{X}, & x^2, & x, & 1, & \epsilon\sqrt{X}/x \\
\sqrt{Y}, & y^2, & y, & 1, & \epsilon\sqrt{Y}/y \\
\sqrt{Z}, & z^2, & z, & 1, & \epsilon\sqrt{Z}/z \\
\sqrt{W}, & w^2, & w, & 1, & \epsilon\sqrt{W}/w \\
\Pi, & a^2, & a, & 1, & 0
\end{vmatrix} = 0,
\]
or, what is the same thing,
\[
\begin{vmatrix}
x^2, & x, & 1, & \epsilon\sqrt{X}/x \\
y^2, & y, & 1, & \epsilon\sqrt{Y}/y \\
z^2, & z, & 1, & \epsilon\sqrt{Z}/z \\
w^2, & w, & 1, & \epsilon\sqrt{W}/w \\
a^2, & a, & 1, & 0
\end{vmatrix} = 0;
\]
viz.
\begin{multline*}
\Pi . x - y . x - z . x - w . y - z . y - w . z - w = -\epsilon \{\sqrt{X}.y - z.y - w.y - a.z - w.z - a.w - a \\
+ \sqrt{Y}.z - w.z - a.z - x.w - a.w - x.a - x \\
+ \sqrt{Z}.w - a.w - x.w - y.a - x.a - y.x - y \\
+ \sqrt{W}.a - x.a - y.a - z.x - y.x - z.y - z\},
\end{multline*}

\vspace{1em}

\begin{center}
\textbf{458}
\end{center}

\vspace{0.5em}

\noindent\textbf{[703}

\vspace{0.5em}

or as it may be written
\begin{multline*}
\Pi . x - z . x - w . y - z . y - w = -\{y - z . y - w . a_2 \sqrt{X} - x - z . x - w . a_1 \sqrt{Y}\} \\
+ \{z - x . z - y . a_3 \sqrt{Z} - z - x . z - y . a_4 \sqrt{W}\},
\end{multline*}
an equation for the determination of $\Pi$.

Consider first the expression which multiplies $\epsilon . a - z . a - w$; this is
\[
= \Im\{y - z . y - w . a_2 \sqrt{X} - x - z . x - w . a_1 \sqrt{Y}\}
\]
we have
\[
BE_{12} = \Im\{\sqrt{b_1 c_1 d_1 e_1 f_1 a_2 c_2 d_2 e_2 f_2} - \sqrt{b_2 c_2 d_2 e_2 f_2 a_1 c_1 d_1 e_1 f_1}\},
\]
and multiplying this by
\[
C_{12} \cdot D_{12} = \sqrt{a_1 c_1 d_1 a_2 c_2 d_2}
\]
we derive
\[
BE_{12} \cdot C_{12} \cdot D_{12} = \Im\{c_1 d_2 \sqrt{X} - c_2 d_1 \sqrt{Y}\},
\]
and similarly two other equations; the system may be written
\begin{align*}
BE_{12} \cdot C_{12} \cdot D_{12} &= \Im\{c_1 d_2 \sqrt{X} - c_2 d_1 \sqrt{Y}\}, \\
CE_{12} \cdot D_{12} \cdot B_{12} &= \Im\{d_1 b_2 \sqrt{X} - d_2 b_1 \sqrt{Y}\}, \\
DE_{12} \cdot B_{12} \cdot C_{12} &= \Im\{b_1 c_2 \sqrt{X} - b_2 c_1 \sqrt{Y}\},
\end{align*}
the suffixes on the left-hand side being always $12$.

The five letters $b$, $c$, $d$, $e$, $f$ enter symmetrically, for $BE$ is a mere abbreviation for the double triad $BEF.ACD$; and the like for $CE$, and $DE$. Any three of the five letters other than $a$; and the remaining two letters $e$, $f$ enter symmetrically into these equations.

Multiplying these equations by
\[
\frac{b - z . b - w}{b - c . b - d},\quad \frac{c - z . c - w}{c - d . c - b},\quad \frac{d - z . d - w}{d - b . d - c}
\]
respectively, and then adding, the right-hand side becomes
\[
= \Im\{y - z . y - w . a_2 \sqrt{X} - x - z . x - w . a_1 \sqrt{Y}\}.
\]
Writing
\[
\frac{b - z . b - w}{b - c . b - d} = \frac{1}{c - d . d - b . b - c} \cdot \text{etc.},
\]
for the left-hand side becomes
\[
= \frac{1}{c - d . d - b . b - c} \{c - d . BE_{12} \cdot C_{12} \cdot D_{12} + d - b . CE_{12} \cdot D_{12} \cdot B_{12} + b - c . DE_{12} \cdot B_{12} \cdot C_{12}\},
\]
which for shortness may be written
\[
= \frac{1}{c - d . d - b . b - c} \sum\{c - d . BE_{12} \cdot C_{12} \cdot D_{12}\},
\]
the summation referring to the three terms obtained by the cyclical interchange of the letters $b$, $c$, $d$.

\vspace{1em}

\begin{center}
\textbf{459}
\end{center}

\vspace{0.5em}

\noindent\textbf{[703}

\vspace{0.5em}

The result thus is
\begin{multline*}
\Im\{y - z . y - w . a_2 \sqrt{X} - x - z . x - w . a_1 \sqrt{Y}\} \\
= \frac{1}{c - d . d - b . b - c} \sum\{c - d . BE_{12} \cdot C_{12} \cdot D_{12}\}.
\end{multline*}

Interchanging $x$, $y$ with $z$, $w$ respectively, we have of course to interchange the suffixes $1$, $2$ and $3$, $4$; we thus find
\begin{multline*}
\Im\{w - x . w - y . a_4 \sqrt{Z} - z - x . z - y . a_3 \sqrt{W}\} \\
= \frac{1}{c - a . a - b . b - c} \sum\{c - d . BE_{34} \cdot C_{34} \cdot D_{34}\}.
\end{multline*}

But $\Pi = \alpha a^2 + \beta a^2 + \gamma a + \delta$, and we hence find the value of $\Pi . x - z . x - w . y - z . y - w$.

The resulting equation divides by $A_{12}^2 \cdot A_{34}^2$: throwing out this factor, we have
\begin{multline*}
\sqrt{a^2 - \xi^2} = A_{56}^{-2} \{x - z . x - w . y - z . y - w\}(c - d . d - b . b - c) A_{12}^{-2} \\
\sum\{c - d . BE_{12} \cdot C_{12} \cdot D_{12}\} + A_{34}^{-2} \sum\{c - d . BE_{34} \cdot C_{34} \cdot D_{34}\},
\end{multline*}
where, as before, the summations refer to the three terms obtained by the cyclical interchange of the letters $b$, $c$, $d$; and the remaining two letters $e$, $f$ enter symmetrically into the formula.

The formula gives thus for $A_{12}$ ten values which are of course equal to each other. Writing for $a$ each letter in succession, we obtain formulae for each of the six single-letter functions $A_{12}$ of $p$ and $q$; and the factor $\sqrt{a^2 - \xi^2}$ is the same in all; and the factor $(x - z . x - w . y - z . y - w)$ is common to all the formulae.

We require further the expressions for the double-letter functions of $p$, $q$. Considering for example the function $BE_{12}$, which is
\[
= \Im\{\sqrt{b_1 c_1 d_1 e_1 f_1 a_2 c_2 d_2 e_2 f_2} - \sqrt{b_2 c_2 d_2 e_2 f_2 a_1 c_1 d_1 e_1 f_1}\}.
\]

\vspace{1em}

\begin{center}
\textbf{460}
\end{center}

\vspace{0.5em}

\noindent\textbf{[703}

\vspace{0.5em}

But the expressions are respectively
\[
\sqrt{a^2 - \xi^2} A_{12} A_{34} A_{56}
\]
and
\[
\frac{1}{\sqrt{a^2 - \xi^2}} A_{12} A_{34} A_{56};
\]
the whole equation thus divides by $A_{12}^2 A_{34}^2$; and then multiplying each side by $\dfrac{\sqrt{a^2 - \xi^2}}{A_{56}}$, throwing out this factor, we find
\[
DE_{12} = \frac{1}{b - c . c - a . a - b} \sum\{c - d . C_{12} \cdot C_{34} \cdot C_{56}\},
\]
in
\[
\sqrt{a^2 - b^2} A_{12} A_{34} A_{56},
\]
which formula if we imagine $\dfrac{\sqrt{a^2 - \xi^2}}{\epsilon}$ each replaced by its value in terms of the $xy$-functions and $zw$-functions, we have an equation of the form
\[
(x - z . x - w . y - z . y - w)\, DE_{12} = \text{a given rational and integral function of the $16$ functions}
\]
where $M$ is a given rational and integral function of the $16$ functions $A_{12}$, $AB_{12}$ and $A_{34}$, $AB_{34}$ of $x$ and $y$ and of $z$ and $w$ respectively.

The factor $(x - z . x - w . y - z . y - w)$ on the left-hand side is retained as being the same factor which enters into the equations for $A_{12}$, etc.: but on the right-hand side $x - z . x - w . y - z . y - w$ should be expressed in terms of the $xy$- and $zw$-functions. This can be done by means of the identity
\[
x - z . x - w . y - z . y - w = \frac{\begin{vmatrix} 1, & x + y, & xy \\ 1, & z + w, & zw \\ 1, & a + b, & ab \end{vmatrix} \cdot \begin{vmatrix} 1, & x + y, & xy \\ 1, & z + w, & zw \\ 1, & a + c, & ac \end{vmatrix}}{(a - b)^2 (a - c)^2},
\]
where the summation refers to the three terms obtained by the cyclical interchange of the letters $a$, $b$, $c$.

The first determinant, multiplied by $a - b$, is in fact
\[
\begin{vmatrix} a - z . a - w, & a - x . a - y \\ b - z . b - w, & b - x . b - y \end{vmatrix}
\]
and the second determinant, multiplied by $a - c$, is
\[
\begin{vmatrix} a - z . a - w, & a - x . a - y \\ c - z . c - w, & c - x . c - y \end{vmatrix},
\]
so that the formula may also be written
\[
x - z . x - w . y - z . y - w = \frac{\begin{vmatrix} a - z . a - w, & a - x . a - y \\ b - z . b - w, & b - x . b - y \end{vmatrix} \cdot \begin{vmatrix} a - z . a - w, & a - x . a - y \\ c - z . c - w, & c - x . c - y \end{vmatrix}}{(a - b)^2 (a - c)^2},
\]
or, what is the same thing, it is
\[
x - z . x - w . y - z . y - w = \frac{\{L + M(a + b) + Nab\}\{L + M(a + c) + Nac\}}{(a - b)^2 (a - c)^2},
\]
where $L$, $M$, $N$ are
\[
= (x + y)zw - (z + w)xy,\quad xy - zw,\quad \text{and}\quad -(x + y) + (z + w);
\]
reduced this is readily verified.

\vspace{1em}

\begin{flushright}
Cambridge, 12th March, 1879.
\end{flushright}

\vspace{1em}

I take this opportunity of remarking that, in the double-letter formulae, the sign of the second term is, not as I have in general written $AB = \Im\{\sqrt{abc_1 d_1 e_1 f_1}\}$, but is $+$, introducing a factor which is a function of $x$ and $y$, the odd and even $\vartheta$-functions are $= \wp\sqrt{a_1 a_2}$, etc.; and this being so, when $x$ and $y$ are interchanged, each single-letter function changes only its sign; and each double-letter function remains unaltered; and thus, since this is a function which on the interchange of $x$, $y$ changes its sign, there is a change of sign as there should be.

\vspace{1em}

\begin{flushright}
Cambridge, 29th July, 1879.
\end{flushright}

\vspace{2em}

% ============================================================
% PAPER 704: A MEMOIR ON THE SINGLE AND DOUBLE THETA-FUNCTIONS
% ============================================================
\begin{center}
\large\bfseries 704.\par\vspace{0.3em}
A MEMOIR ON THE SINGLE AND DOUBLE THETA-FUNCTIONS.\par\vspace{0.5em}
\normalsize\normalfont
[From the Philosophical Transactions of the Royal Society of London, vol.~171, Part~III.\ (1880), pp.~897---1002.\ Received November~14,--- Read November~28, 1879.]
\end{center}

\vspace{1em}

The Theta-Functions, although arising historically from the Elliptic Functions, may be considered as in order of simplicity preceding these, and connecting themselves directly with the exponential function $(e^x$ or$)\ \exp. x$; viz.\ they may be defined each of them as a sum of a series of exponentials, singly infinite in the case of the single functions, doubly infinite in the case of the double functions, and so on. The number of the single functions is $= 4$; the number of the double functions is $= 16$ $(4^2 =)$; and the quotients of these, or say three of them each divided by the fourth, are the elliptic functions $\sn$, $\cn$, $\dn$; and the quotients of these, or say fifteen of them each divided by the sixteenth, are the hyper-elliptic functions of two arguments depending on the square root of a sextic function. Generally, the number of the $p$-tuple theta-functions is $= 4^p$; and the quotients of these, or say all but one of them each divided by the remaining function, are the Abelian functions of $p$ arguments depending on the irrationality $y$ defined by the equation $F(x, y) = 0$ of a curve of deficiency $p$.

If, instead of the ratios of the functions, we consider the ratios of the functions connecting the ratios of $(4p - 1)$ dimensions, then we have the sixteen coordinates of a point on a quadri-quadric two-fold locus in $15$-dimensional space, the deficiency of this two-fold locus being of course $= 2$.

\vspace{1em}

\begin{center}
\textbf{463}
\end{center}

\vspace{0.5em}

\noindent\textbf{[704}

\vspace{0.5em}

The Memoir relates only to the single and double functions, and the title has been given accordingly. The investigations contained in it extend to the single functions referred to in the First Part; but in the main and with a view to simplicity of notation the investigations are applicable to the general case of the double functions; but in the present Memoir, although the course of the single theory is in fact more complex than that of the double theory; and on the other hand we need the single theory as a guide for the course of the double theory.

I accordingly stop to point out that the First Part of the Memoir relates to the single functions; the Second and the Third Parts to the double functions respectively; and I consider in the first place the course of the single theory; and then that of the double theory. The paragraphs of the Memoir are numbered consecutively.

The general definition adopted for the theta-functions differs somewhat from that ordinarily used. The earlier memoirs on the double theta-functions are the well-known ones

\noindent Rosenhain, ``M\'{e}moire sur les fonctions de deux variables \`{a} quatre p\'{e}riodes, qui sont les inverses des int\'{e}grales ultra-elliptiques de la premi\`{e}re classe.'' M\'{e}m.\ Savans \'{E}trang., t.~XI.\ (1851), pp.~361---468.

\noindent Gopel, ``Theoriae transcendentium Abelianarum primi ordinis adumbratio levis,'' Crelle, t.~XXXV.\ (1847), pp.~277---312.

\noindent My first paper---Cayley, ``On the Double $\vartheta$-Functions in connexion with a 16-nodal Surface,'' Crelle, t.~LXXXIII.\ (1877), pp.~210---219, [662]---was founded directly upon these, and was immediately followed by Dr Borchardt's paper,

\noindent Borchardt, ``Ueber die Darstellung der Kummer'schen Fl\"{a}che vierter Ordnung mit sechzehn Knotenpunkten durch die Gopelsche biquadratische Relation zwischen vier Thetafunctionen mit zwei Variabeln,'' Ditto, pp.~234---244.

My other later papers, [663, 664, 665, 697, 703], are contained in the same Journal.

\vspace{1em}

\begin{center}
\textsc{First Part.---Introductory.}
\end{center}

\vspace{0.5em}

\begin{center}
\textit{Definition of the theta-functions.}
\end{center}

\vspace{0.5em}

\noindent\textbf{1.}\quad
The $p$-tuple theta-functions depend upon $\frac{1}{2}p(p+1)$ parameters which are the coefficients of a quadric function of $p$ ultimately disappearing integers, upon $p$ arguments, and upon $2p$ characters, each $= 0$ or $1$, which form the characteristic of the $4^p$ functions; but it will be sufficient to write down the formulae in the case $p = 2$.

As already mentioned, the adopted definition differs somewhat from that ordinarily used. I use, as will be seen, a quadric function $J(a, h, b\cbtw m, n)^2$ with even or odd values; and I write the other term $\frac{1}{2}\pi i(mu + nv)$, instead of $mu + nv$; this comes to affecting the arguments $u$, $v$ with a factor $\pi i$, so that the quarter-periods (instead of being $\pi i$) are made to be $= 1$.

\noindent\textbf{2.}\quad
We write
\[
\begin{pmatrix} m, & n \\ u, & v \end{pmatrix} = J(a, h, b\cbtw m, n)^2 + \tfrac{1}{2}\pi i(mu + nv),
\]
and in like manner
\[
\begin{pmatrix} m + \alpha, & n + \beta \\ u + \gamma, & v + \delta \end{pmatrix} = J(a, h, b\cbtw m + \alpha, n + \beta)^2 + \tfrac{1}{2}\pi i\{(m + \alpha)(u + \gamma) + (n + \beta)(v + \delta)\},
\]
and prefixing to either of these the functional symbol $\exp.$ we have the exponential of the function in question, that is, $e$ with the function as an exponent.

We then write, as the definition of the double theta-functions,
\[
\vartheta\begin{pmatrix} \alpha, & \beta \\ \gamma, & \delta \end{pmatrix}(u, v) = \Sigma\, \exp.\begin{pmatrix} m + \alpha, & n + \beta \\ u + \gamma, & v + \delta \end{pmatrix},
\]
where the summation extends to all positive and negative even integer values (zero included) of $m$ and $n$ respectively; $\alpha$, $\beta$, $\gamma$, $\delta$ might denote any quantities whatever, but for the theta-functions they are regarded as denoting positive or negative integers; this being so, it will appear that the only effect of altering each or any of them by an even integer is to reverse (it may be) the sign of the function; and the distinct functions are consequently the $(4^2 =)\ 16$ functions obtained by giving to each of the quantities $\alpha$, $\beta$, $\gamma$, $\delta$ the two values $0$ and $1$ successively.

\noindent\textbf{3.}\quad
We thus have the double theta-functions, depending on the parameters $(a, h, b)$ which determine the quadric function $(a, h, b\cbtw m, n)^2$ of the disappearing even integers $(m, n)$, and on the two arguments $(u, v)$: in the symbol $\begin{pmatrix} \alpha, & \beta \\ \gamma, & \delta \end{pmatrix}$, which is called the characteristic, the characters $\alpha$, $\beta$, $\gamma$, $\delta$ are each of them $= 0$ or $1$; and we thus have the $16$ functions.

The parameters $(a, h, b)$ may be real or imaginary, but they must be such that reducing each of them to its real part the resulting function $(a, h, b\cbtw m, n)^2$ is invariable in its sign, and negative for all real values of $m$ and $n$: this is, in fact, the condition for the convergency of the series which give the values of the theta-functions.

\noindent\textbf{4.}\quad
The characteristic $\begin{pmatrix} \alpha, & \beta \\ \gamma, & \delta \end{pmatrix}$ is said to be even or odd according as the sum $\alpha\gamma + \beta\delta$ is even or odd.

\vspace{1em}

\begin{center}
\textit{Allied functions.}
\end{center}

\vspace{0.5em}

\noindent\textbf{5.}\quad
As already remarked, the definition of
\[
\vartheta\begin{pmatrix} \alpha, & \beta \\ \gamma, & \delta \end{pmatrix}(u, v)
\]
is not restricted to the case where the characters $\alpha$, $\beta$, $\gamma$, $\delta$ are integers; and the functions thus obtained are not regarded as theta-functions, and there is in this no occasion to consider functions of the form where they are not integers; but the expression theta-function will consequently not extend to include them.

\vspace{1em}

\begin{center}
\textit{Properties of the Theta-Functions: Various sub-headings.}
\end{center}

\vspace{0.5em}

\begin{center}
\textit{Even-integer alteration of characters.}
\end{center}

\vspace{0.5em}

\noindent\textbf{6.}\quad
If $x$, $y$ be even integers, then $m$, $n$ having the several even integer values from $-\infty$ to $+\infty$ respectively, it is obvious that $m + \alpha$, $n + \beta$ and $m + \alpha + 2x$, $n + \beta + 2y$ will have the same series of values; and it thence follows that
\[
\vartheta\begin{pmatrix} \alpha + 2x, & \beta + 2y \\ \gamma, & \delta \end{pmatrix}(u, v) = \vartheta\begin{pmatrix} \alpha, & \beta \\ \gamma, & \delta \end{pmatrix}(u, v).
\]
Similarly if $z$, $w$ are integers, then in the function
\[
\vartheta\begin{pmatrix} \alpha, & \beta \\ \gamma + 2z, & \delta + 2w \end{pmatrix}(u, v)
\]
the argument of the exponential function contains the term
\[
\tfrac{1}{2}\pi i\{(m + \alpha)(u + \gamma + 2z) + (n + \beta)(v + \delta + 2w)\}
\]
this differs from its original value by
\[
\tfrac{1}{2}\pi i\{(m + \alpha)2z + (n + \beta)2w\} = \pi i(mz + nw) + \pi i(\alpha z + \beta w),
\]
and then, $m$ and $n$ being even integers, $mz + nw$ is also an even integer, and the term $\pi i(mz + nw)$ does not affect the value of the exponential: we thus introduce into each term of the series the factor $\exp.\pi i(\alpha z + \beta w)$, which is, in fact, $= (-)^{\alpha z + \beta w}$; and we consequently have
\[
\vartheta\begin{pmatrix} \alpha, & \beta \\ \gamma + 2z, & \delta + 2w \end{pmatrix}(u, v) = (-)^{\alpha z + \beta w}\vartheta\begin{pmatrix} \alpha, & \beta \\ \gamma, & \delta \end{pmatrix}(u, v);
\]
or, uniting the two results,
\[
\vartheta\begin{pmatrix} \alpha + 2x, & \beta + 2y \\ \gamma + 2z, & \delta + 2w \end{pmatrix}(u, v) = (-)^{\alpha z + \beta w}\vartheta\begin{pmatrix} \alpha, & \beta \\ \gamma, & \delta \end{pmatrix}(u, v).
\]
This sustains the before-mentioned conclusion that the only distinct functions are the $16$ functions obtained by giving to the characters $\alpha$, $\beta$, $\gamma$, $\delta$ the values $0$ and $1$ respectively.

\vspace{1em}

\begin{center}
\textit{Odd-integer alteration of characters.}
\end{center}

\vspace{0.5em}

\noindent\textbf{7.}\quad
The effect is obviously to interchange the different functions.

\vspace{1em}

\begin{center}
\textit{Even and odd functions.}
\end{center}

\vspace{0.5em}

\noindent\textbf{8.}\quad
It is clear that $-m - \alpha$, $-n - \beta$ have precisely the same series of values as $m + \alpha$, $n + \beta$ respectively; hence considering the function:
\[
\vartheta\begin{pmatrix} \alpha, & \beta \\ \gamma, & \delta \end{pmatrix}(-u, -v)
\]
the linear term in the argument of the exponential
\[
\tfrac{1}{2}\pi i\{(-m - \alpha)(-u + \gamma) + (-n - \beta)(-v + \delta)\}
\]
is
\[
\tfrac{1}{2}\pi i\{(m + \alpha)(u - \gamma) + (n + \beta)(v - \delta)\}
\]
which differs from its original value
\[
\tfrac{1}{2}\pi i\{(m + \alpha)(u + \gamma) + (n + \beta)(v + \delta)\}
\]
by $-\pi i\{(m + \alpha)\gamma + (n + \beta)\delta\}$, that is $-\pi i(m\gamma + n\delta) - \pi i(\alpha\gamma + \beta\delta)$; the term $-\pi i(m\gamma + n\delta)$, where $m$, $\gamma$ and $n$, $\delta$ are integers, does not affect the value of the exponential; and the term $-\pi i(\alpha\gamma + \beta\delta)$ introduces into each term of the series the factor $(-)^{\alpha\gamma + \beta\delta}$, that is, $+1$ or $-1$ according as the characteristic is even or odd. Hence for an even characteristic the function is even, and for an odd characteristic the function is odd; that is,
\[
\vartheta\begin{pmatrix} \alpha, & \beta \\ \gamma, & \delta \end{pmatrix}(-u, -v) = \pm\vartheta\begin{pmatrix} \alpha, & \beta \\ \gamma, & \delta \end{pmatrix}(u, v),
\]
the upper or lower sign according as the characteristic is even or odd.

\vspace{1em}

\begin{center}
\textit{Change of signs of arguments.}
\end{center}

\vspace{0.5em}

\noindent\textbf{9.}\quad
Changing the signs of the arguments, or changing the signs of the characters $\gamma$, $\delta$, it is clear that we have
\[
\vartheta\begin{pmatrix} \alpha, & \beta \\ \gamma, & \delta \end{pmatrix}(-u, -v) = \vartheta\begin{pmatrix} \alpha, & \beta \\ -\gamma, & -\delta \end{pmatrix}(u, v),
\]
and moreover, altering $\gamma$, $\delta$ each by an even integer, the sign is not altered; hence the signs of $\gamma$, $\delta$ may be simultaneously changed; that is,
\[
\vartheta\begin{pmatrix} \alpha, & \beta \\ \gamma, & \delta \end{pmatrix}(u, v) = \vartheta\begin{pmatrix} \alpha, & \beta \\ -\gamma, & -\delta \end{pmatrix}(-u, -v).
\]

\vspace{1em}

\begin{center}
\textit{Interchange of characters and arguments.}
\end{center}

\vspace{0.5em}

\noindent\textbf{10.}\quad
Writing in the original formula $m$, $n$ for $n$, $m$ respectively, and at the same time interchanging $u$ with $v$ and $\alpha$ with $\beta$, and $\gamma$ with $\delta$, the function is unaltered; that is,
\[
\vartheta\begin{pmatrix} \alpha, & \beta \\ \gamma, & \delta \end{pmatrix}(u, v) = \vartheta\begin{pmatrix} \beta, & \alpha \\ \delta, & \gamma \end{pmatrix}(v, u).
\]

\vspace{1em}

\begin{center}
\textit{Zero arguments.}
\end{center}

\vspace{0.5em}

\noindent\textbf{11.}\quad
The function $\vartheta\begin{pmatrix} \alpha, & \beta \\ \gamma, & \delta \end{pmatrix}(0, 0)$, say for shortness $\vartheta\begin{pmatrix} \alpha, & \beta \\ \gamma, & \delta \end{pmatrix}$, vanishes if the characteristic be odd; but if the characteristic be even, it is in general different from zero. We write
\[
A = \vartheta\begin{pmatrix} 0, & 0 \\ 0, & 0 \end{pmatrix},\quad
B = \vartheta\begin{pmatrix} 0, & 0 \\ 1, & 0 \end{pmatrix},\quad
C = \vartheta\begin{pmatrix} 0, & 0 \\ 0, & 1 \end{pmatrix},\quad
D = \vartheta\begin{pmatrix} 0, & 0 \\ 1, & 1 \end{pmatrix}.
\]
And similarly in the general case we denote the $2^{p-1}(2^p + 1)$ even theta-functions by $A$, $B$, etc., and the $2^{p-1}(2^p - 1)$ odd theta-functions by the same letters.

\vspace{1em}

\begin{center}
\textit{Product-theorem.}
\end{center}

\vspace{0.5em}

\noindent\textbf{12.}\quad
The product of any two theta-functions of the same arguments $(u, v)$ and $(u', v')$ respectively is equal to the sum of four products of theta-functions of the arguments $(u + u', v + v')$ and $(u - u', v - v')$ respectively. In fact, taking the product
\[
\vartheta\begin{pmatrix} \alpha, & \beta \\ \gamma, & \delta \end{pmatrix}(u, v)
\vartheta\begin{pmatrix} \alpha', & \beta' \\ \gamma', & \delta' \end{pmatrix}(u', v')
\]
the exponential of any term hereof is
\[
\exp.\left[\begin{pmatrix} m + \alpha, & n + \beta \\ u + \gamma, & v + \delta \end{pmatrix} + \begin{pmatrix} m' + \alpha', & n' + \beta' \\ u' + \gamma', & v' + \delta' \end{pmatrix}\right],
\]
or writing $M = m + m'$, $N = n + n'$, $m = \mu + M$, $m' = \mu' + M$, etc., this is
\begin{multline*}
\exp.\left[\begin{pmatrix} M + \tfrac{1}{2}(\alpha + \alpha'), & N + \tfrac{1}{2}(\beta + \beta') \\ u + u' + \tfrac{1}{2}(\gamma + \gamma'), & v + v' + \tfrac{1}{2}(\delta + \delta') \end{pmatrix}\right. \\
+ \left.\begin{pmatrix} \mu + \tfrac{1}{2}(\alpha - \alpha'), & \nu + \tfrac{1}{2}(\beta - \beta') \\ u - u' + \tfrac{1}{2}(\gamma - \gamma'), & v - v' + \tfrac{1}{2}(\delta - \delta') \end{pmatrix}\right],
\end{multline*}
where observe that the even integer values of $m$, $m'$ give for $M$, $\mu$ even integer values, but these are not independent; we have in fact $M + \mu = m + m'$, $M - \mu = m - m'$, which are both even, that is, $M$ and $\mu$ are of the same parity; the terms are consequently $M$ even, $\mu$ even, or else $M$ odd, $\mu$ odd; the pairs of values are consequently $(M, \mu) = (\text{even}, \text{even})$, or $(\text{odd}, \text{odd})$. And similarly $(N, \nu) = (\text{even}, \text{even})$, or $(\text{odd}, \text{odd})$. The pair of conditions may be satisfied in four different ways; and for each of these the sum extends to all positive and negative even integer values of $M$, $N$, $\mu$, $\nu$; that is, it extends to all positive and negative even integer values of $m$, $n$, $m'$, $n'$.

Hence arranging the terms according to the four cases respectively, and observing that in the four cases respectively the terms of the product contain the four factors
\[
(-)^{\frac{1}{2}(\alpha - \alpha')(\gamma - \gamma') + \frac{1}{2}(\beta - \beta')(\delta - \delta')},\quad
(-)^{\frac{1}{2}(\alpha - \alpha')(\gamma - \gamma') + \frac{1}{2}(\beta - \beta')(\delta - \delta') + 1},
\]
etc., we obtain the theorem.

\noindent\textbf{13.}\quad
The product of two functions where the two characteristics are
\[
\begin{pmatrix} \alpha, & \beta \\ \gamma, & \delta \end{pmatrix}\quad \text{and}\quad
\begin{pmatrix} \alpha', & \beta' \\ \gamma', & \delta' \end{pmatrix}
\]
respectively, is equal to the sum of the four products which correspond to the four characteristics
\[
\begin{pmatrix} \tfrac{1}{2}(\alpha + \alpha'), & \tfrac{1}{2}(\beta + \beta') \\ \tfrac{1}{2}(\gamma + \gamma'), & \tfrac{1}{2}(\delta + \delta') \end{pmatrix}
\begin{pmatrix} \tfrac{1}{2}(\alpha - \alpha'), & \tfrac{1}{2}(\beta - \beta') \\ \tfrac{1}{2}(\gamma - \gamma'), & \tfrac{1}{2}(\delta - \delta') \end{pmatrix}
\]
with the proper signs. In particular, when the two characteristics are the same, we have the square of a single function equal to the sum of four products.

\noindent\textbf{14.}\quad
The cases most required are those in which each of the characteristics $\begin{pmatrix} \alpha, & \beta \\ \gamma, & \delta \end{pmatrix}$, $\begin{pmatrix} \alpha', & \beta' \\ \gamma', & \delta' \end{pmatrix}$ has its four characters each $= 0$ or $1$; and consequently the four characteristics in the reduced product have their characters $= 0$ or $\tfrac{1}{2}$. It is easy to see that if in any one of these we alter any character by an integer (that is, here by $0$ or by $1$), the effect is to interchange the functions; hence it is only necessary to write down a set of formulae for the cases where the characteristics in the reduced product have their characters each $= 0$ or $\tfrac{1}{2}$, but subject to no other condition; the signs are to be determined by the formula.

\vspace{1em}

\begin{center}
\textit{Resume of the ulterior theory of the single functions.}
\end{center}

\vspace{0.5em}

\noindent\textbf{15.}\quad
For the single functions, I attend here to the full theory in the sequel. But attending at present to the single functions, I write down four of the $16$ equations, viz.\ these are

\begin{center}
\begin{tabular}{c@{\qquad}c}
$\vartheta\begin{pmatrix} 0 \\ 0 \end{pmatrix}(u + u') \cdot \vartheta\begin{pmatrix} 0 \\ 0 \end{pmatrix}(u - u')$ & $= XX' + YY'$, \\
$\vartheta\begin{pmatrix} 1 \\ 0 \end{pmatrix}(u + u') \cdot \vartheta\begin{pmatrix} 1 \\ 0 \end{pmatrix}(u - u')$ & $= YX' + XY'$, \\
$\vartheta\begin{pmatrix} 0 \\ 1 \end{pmatrix}(u + u') \cdot \vartheta\begin{pmatrix} 0 \\ 1 \end{pmatrix}(u - u')$ & $= XX' - YY'$, \\
$\vartheta\begin{pmatrix} 1 \\ 1 \end{pmatrix}(u + u') \cdot \vartheta\begin{pmatrix} 1 \\ 1 \end{pmatrix}(u - u')$ & $= -YX' + XY'$;
\end{tabular}
\end{center}

\noindent where $X$, $Y$ denote $\vartheta\begin{pmatrix} 0 \\ 0 \end{pmatrix}(2u)$, $\vartheta\begin{pmatrix} 1 \\ 0 \end{pmatrix}(2u)$ respectively, and $X'$, $Y'$ the same functions of $2u'$ respectively.

\vspace{1em}

\textit{Product-theorem for the single theta-functions.}

\noindent In the other $12$ equations we have on the left-hand the product of different theta-functions of $u + u'$, $u - u'$ respectively, and on the right-hand expressions involving other functions, $X_1$, $Y_1$, $X_1'$, $Y_1'$, etc., of $2u$ and $2u'$ respectively.

Writing $u' = 0$, we have on the left-hand, squares or products of the theta-functions of $u$, and on the right-hand expressions containing functions of $2u$; in particular, the above equations show that the squares of the four theta-functions are equal to linear functions of $X$, $Y$; that is, there exist between the squared functions two linear relations.

\vspace{1em}

\begin{center}
\textbf{471}
\end{center}

\vspace{0.5em}

\noindent\textbf{[704}

\vspace{0.5em}

\noindent\textbf{16.}\quad
By introducing a variable argument $x$, the four squared functions may be taken to be proportional to linear functions of $x$:
\[
\begin{matrix}
A^2 u & : & B^2 u & : & C^2 u & : & D^2 u \\
\mathfrak{A}(a - x) & : & \mathfrak{A}(b - x) & : & \mathfrak{C}(c - x) & : & \mathfrak{D}(d - x),
\end{matrix}
\]
where $a$, $b$, $c$, $d$ are constants.

This suggests a new notation; we write
\begin{align*}
A^2 u &= \omega^2 \mathfrak{A}(a - x), \\
B^2 u &= \omega^2 \mathfrak{A}(b - x), \\
C^2 u &= \omega^2 \mathfrak{C}(c - x), \\
D^2 u &= \omega^2 \mathfrak{D}(d - x);
\end{align*}
and the result just mentioned then is
\[
A^2 u : B^2 u : C^2 u : D^2 u = \mathfrak{A}(a - x) : \mathfrak{A}(b - x) : \mathfrak{C}(c - x) : \mathfrak{D}(d - x),
\]
which expresses that the four functions are the coordinates of a point on a quadri-quadric curve in ordinary space.

The remaining $12$ of the $16$ equations then contain on the left-hand products such as
\[
A(u + u') \cdot B(u - u');
\]
and by suitably combining them we obtain equations such as
\[
\frac{A(u + u') \cdot B(u - u')}{C(u + u') \cdot D(u - u') + D(u + u') \cdot C(u - u')} = \text{function of } u',
\]
where for brevity the arguments are as written above; viz.\ the numerator of the fraction is
\[
B(u + u') \cdot A(u - u') - A(u + u') \cdot B(u - u'),
\]
and its denominator is
\[
C(u + u') \cdot D(u - u') + D(u + u') \cdot C(u - u').
\]
Admitting the form of the equation, the value of the function of $u'$ is at once found by writing in the equation $u = 0$; it is, as it ought to be, a function vanishing for $u' = 0$.

\noindent\textbf{17.}\quad
Take $u$ indefinitely small; the resulting equation appears thus:
\[
\frac{Au \cdot Bu' - Bu \cdot Au'}{Cu \cdot Du' + Du \cdot Cu'} = \text{const.},
\]
where $Au$, $Bu$ are the derived functions, or dividing each side by $u'$, the differential coefficients; or, what is the same thing, that the differential coefficient of the same combination is a constant in regard to $u$.

\noindent\textbf{18.}\quad
It appears thus that $\dfrac{Au \cdot Bu}{Cu \cdot Du}$ is a constant multiple of the product of the two quotient-functions $p$ and $q$. And then substituting for the several quotient-functions their values in terms of $x$, we obtain a differential relation between $x$, $u$ and $v$; viz.
\[
\frac{dx}{\sqrt{a - x . b - x . c - x . d - x}} = M\, du,
\]
and it is thus that the quotient-functions are as obtained in the sequel; viz.
\[
\sn Ku = \sqrt{\frac{FAu}{Bu}} = Du \div Cu,
\]
\[
\cn Ku = \sqrt{\frac{GAu}{Bu}} = Bu \div Cu,
\]
\[
\dn Ku = \sqrt{\frac{HAu}{Bu}};
\]
in fact, the elliptic-functions:
\[
\sn u,\quad \cn u,\quad \dn u
\]
and we thus identify the functions $Au$, $Bu$, $Cu$, $Du$ with the functions of Jacobi.

If in the above-mentioned four equations we write first $u = 0$, and then $u' = 0$, by means of the results eliminate from the original equations the quantities $X$, $Y$, $X'$, $Y'$ which occur therein, we obtain expressions for the four products such as $A(u + u') \cdot A(u - u')$.

\noindent\textbf{19.}\quad
One of these equations is
\[
C(u + u') \cdot C(u - u') = Cu \cdot Cu' - Du \cdot Du',
\]
and taking herein $u'$ indefinitely small, we obtain
\[
\frac{Cu \cdot C''u - (C'u)^2}{Cu^2} = \frac{d^2}{du^2}\log Cu,
\]
where the left-hand side is given in fact in terms of the theta-function; and taking the exponential of each side twice, we obtain $Cu$ in terms of the function.

This, in fact, corresponds to Jacobi's equation
\[
\frac{d^2}{du^2}\log \vartheta\begin{pmatrix} 0 \\ 1 \end{pmatrix}u = -k^2 \sn^2 u + \text{const.}
\]

\noindent\textbf{20.}\quad
Which contains the double $\xi$; hence, integrating
\[
\int \sn^2 u\, du,
\]
an exponential as the integral of a squared quotient-function.

From the same equation
\[
C(u) \cdot C(u + u') \cdot C(u - u') = Cu \cdot C^2 u' - B^2 u \cdot D^2 u',
\]
differentiating logarithmically in regard to $u'$ and integrating in regard to $u$, we obtain an equation containing on the left-hand side a term $\log \dfrac{B'u}{Au}$, and on the right-hand side an integral in regard to $u$ of $u'$:
\[
\int_0^u \frac{B'u}{Au}\, du = u' \cdot \text{const.} + \log \vartheta\begin{pmatrix} 0 \\ 1 \end{pmatrix}(u - a);
\]
in fact, this corresponds to Jacobi's equation
\[
\int_0^u \sn u \cn u \, du = \sn' u \dn' u.
\]

\noindent\textbf{21.}\quad
It may further be noticed that if, in the other three equations of the system, we introduce the corresponding quantity $\xi$ in place of $u'$, then the integral in question is in place of an expression such as
\[
\int \frac{dx}{\sqrt{a - x . b - x . c - x . d - x}},
\]
where $T$ is $= -\xi$, or is $= $ any one of three forms such as
\[
\begin{vmatrix}
1, & x, & x^2 \\
1, & a + b, & ab \\
1, & c + d, & cd
\end{vmatrix}.
\]

\vspace{1em}

\begin{center}
\textbf{472}
\end{center}

\vspace{0.5em}

\noindent\textbf{[704}

\vspace{0.5em}

\begin{center}
\textit{Resume of the ulterior theory of the double functions.}
\end{center}

\vspace{0.5em}

\noindent\textbf{23.}\quad
As regards the ulterior theory of the double functions, it is to be observed that we have not only the $16$ equations leading to linear relations between binary products of the $16$ different functions. We have thus between the $16$ functions a system of quadric relations, which in fact determine the ratios of the $16$ functions in terms of two variable parameters $x$, $y$. (The $16$ functions are thus the coordinates of a point on a quadri-quadric two-fold locus in $15$-dimensional space.) The forms depend upon six constants, $a$, $b$, $c$, $d$, $e$, $f$:
\[
\sqrt{A} = \sqrt{ah}, \quad \sqrt{a - x . a - y},
\]
\[
\sqrt{Ab} = \sqrt{a - x . b - x . c - x . d - x . e - x . f - x . a - y . b - y . c - y . d - y . e - y . f - y}
\]
(observe that in the symbols $\sqrt{ab}$, etc., the letter $f$ always accompanies the two expressed letters; or, what is the same thing, the duad $ab$ is really an abbreviation for the double triad $abf.cde$): this suggests that the functions should be represented by the single and double letter notation $A(u, v)$, \dots, $AB(u, v)$, \dots; viz.\ if for shortness the arguments are omitted, then we have
\[
A,\quad B,\quad C,\quad D,\quad E,\quad F,\quad AB,\quad AC,\quad AD,\quad AE,\quad AF,\quad BC,\quad BD,\quad BE,\quad BF,\quad CD,\quad CE,\quad CF,\quad DE
\]
proportional to determinate constant multiples of
\[
\sqrt{a},\ \dots,\ \sqrt{ab},\ \dots,
\]
of $x$ and $y$; and the squared relations between the $16$ functions also lead to writing for shortness
\[
\frac{\sqrt{a - x . b - x . c - x . d - x . e - x . f - x . a - y . b - y . c - y . d - y . e - y . f - y}}{x - y}
\]

\noindent\textbf{24.}\quad
It is interesting to notice why in the expressions for $\sqrt{ab}$, etc., the effect of the interchange of $x$, $y$ is, in fact, to change $(u, v)$ into $(-u, -v)$; consequently to change the sign of the odd functions, and to leave unaltered those of the even functions: the interchange does in fact leave $\sqrt{a}$, etc., unaltered, while it changes $\sqrt{ab}$, etc., into $-\sqrt{ab}$, etc.; and thus, since only the ratios are attended to, there is a change of sign as there should be.

\noindent\textbf{25.}\quad
The equations of the product-theorem lead to expressions for
\[
\begin{matrix} u + u' & u - u' \\ A \cdot B - B \cdot A, \end{matrix}
\]
where the arguments, written above, are used to denote the two arguments, viz.\ $u + u'$ to denote $(u + u', v + v')$ and $u - u'$ to denote $(u - u', v - v')$; and where the letters $A$, $B$ denote each or either of them a single or double letter. These expressions are found in terms of the functions of $(u, v)$ and of $(u', v')$:

taking any such expression in $u'$, $v'$ each of them indefinitely small, but with their ratio arbitrary, we obtain the value of
\[
A \cdot dB - B \cdot dA,
\]
(viz.\ $w$ here stands for the two arguments $(u, v)$, and $\mathfrak{d}$ denotes total differentiation $dA = du \cdot \dfrac{d}{du}A(u, v) + dv \cdot \dfrac{d}{dv}A(u, v)$), as a quadric function of the functions of $(u, v)$; or dividing by $A^2$, the form of $\mathfrak{d} \cdot \dfrac{B}{A}$, that we should equal the differentials of the quotient-functions themselves.

Substituting for the quotient-functions their values in terms of $x$, $y$, we have the differential relations between $dx$, $dy$, $du$, $dv$, viz.\ putting for shortness
\[
X = a - x . b - x . c - x . d - x . e - x . f - x,
\]
and
\[
Y = a - y . b - y . c - y . d - y . e - y . f - y,
\]
these are of the form
\[
\frac{dx}{\sqrt{X}} = \frac{x\,dx}{\sqrt{X}},\quad \frac{dy}{\sqrt{Y}} = \frac{y\,dy}{\sqrt{Y}},
\]
each of them equal to a linear function of $du$ and $dv$; and in fact there are $16$ equations: so that the $16$ functions thus give the hyperelliptic addition-theorem for the quotient-functions belonging to the integrals
\[
\int \frac{dx}{\sqrt{X}},\quad \int \frac{x\,dx}{\sqrt{X}}
\]
in accordance with the theory of these integrals.

\vspace{1em}

\begin{center}
\textbf{473}
\end{center}

\vspace{0.5em}

\noindent\textbf{[704}

\vspace{0.5em}

\noindent\textbf{26.}\quad
The first of the product-theorem, putting therein first $u = 0$, then $v' = 0$, and using the results to eliminate the functions on the right-hand side, give expressions for $A(u + u', v + v') \cdot B(u - u', v - v')$, etc., in terms of the functions of $(u, v)$ and $(u', v')$: that is, they give an addition-with-subtraction theorem for the double theta-functions. And we have thence also consequences analogous to those which present themselves in the theory of the single functions.

\vspace{1em}

\begin{center}
\textit{Remark as to notation.}
\end{center}

\vspace{0.5em}

\noindent\textbf{27.}\quad
I remark, as regards the single theta-functions, that the characteristics $\vartheta\begin{pmatrix} 0 \\ 0 \end{pmatrix}$, $\vartheta\begin{pmatrix} 1 \\ 0 \end{pmatrix}$, $\vartheta\begin{pmatrix} 0 \\ 1 \end{pmatrix}$, $\vartheta\begin{pmatrix} 1 \\ 1 \end{pmatrix}$ might for shortness be represented by a series of current numbers $0$, $1$, $2$, $3$; and the functions be accordingly called $\vartheta_0 u$, $\vartheta_1 u$, $\vartheta_2 u$, $\vartheta_3 u$; but that, instead of this, I prefer to use throughout the before-mentioned functional symbols $A$, $B$, $C$, $D$.

As regards the double functions, I do, however, denote the characteristics
\[
\begin{matrix}
\begin{pmatrix} 0, & 0 \\ 0, & 0 \end{pmatrix}, &
\begin{pmatrix} 0, & 0 \\ 1, & 0 \end{pmatrix}, &
\begin{pmatrix} 0, & 0 \\ 0, & 1 \end{pmatrix}, &
\begin{pmatrix} 0, & 0 \\ 1, & 1 \end{pmatrix}, \\[1em]
\begin{pmatrix} 1, & 0 \\ 0, & 0 \end{pmatrix}, &
\begin{pmatrix} 1, & 0 \\ 1, & 0 \end{pmatrix}, &
\begin{pmatrix} 1, & 0 \\ 0, & 1 \end{pmatrix}, &
\begin{pmatrix} 1, & 0 \\ 1, & 1 \end{pmatrix}, \\[1em]
\begin{pmatrix} 0, & 1 \\ 0, & 0 \end{pmatrix}, &
\begin{pmatrix} 0, & 1 \\ 1, & 0 \end{pmatrix}, &
\begin{pmatrix} 0, & 1 \\ 0, & 1 \end{pmatrix}, &
\begin{pmatrix} 0, & 1 \\ 1, & 1 \end{pmatrix}, \\[1em]
\begin{pmatrix} 1, & 1 \\ 0, & 0 \end{pmatrix}, &
\begin{pmatrix} 1, & 1 \\ 1, & 0 \end{pmatrix}, &
\begin{pmatrix} 1, & 1 \\ 0, & 1 \end{pmatrix}, &
\begin{pmatrix} 1, & 1 \\ 1, & 1 \end{pmatrix},
\end{matrix}
\]
by the numbers $1$, $2$, $3$, $4$, $5$, $6$, $7$, $8$, $9$, $10$, $11$, $12$, $13$, $14$, $15$, $16$ respectively; or say the characteristics are represented by the numbers $1$ to $16$, and the corresponding functions are $A$, $B$, $C$, $D$, $E$, $F$, $G$, $H$, $I$, $J$, $K$, $L$, $M$, $N$, $O$, $P$.

\vspace{1em}

\begin{center}
\textbf{474}
\end{center}

\vspace{0.5em}

\noindent\textbf{[704}

\vspace{0.5em}

\noindent\textbf{28.}\quad
But the formulae will be written in a more convenient form by using for the characteristics a binary notation; viz.\ writing $\alpha$, $\beta$, $\gamma$, $\delta$ each $= 0$ or $1$, I represent the characteristic by the symbol
\[
\begin{pmatrix} \alpha, & \beta \\ \gamma, & \delta \end{pmatrix}
\]
and the corresponding function by
\[
\vartheta\begin{pmatrix} \alpha, & \beta \\ \gamma, & \delta \end{pmatrix}(u, v)\quad \text{or}\quad \vartheta\begin{pmatrix} \alpha, & \beta \\ \gamma, & \delta \end{pmatrix}.
\]

\vspace{1em}

\begin{center}
\textit{Second Part.---The Single Theta-Functions.}
\end{center}

\vspace{0.5em}

\begin{center}
\textit{Definition and fundamental properties.}
\end{center}

\vspace{0.5em}

\noindent\textbf{29.}\quad
The single theta-functions are defined by the equation
\[
\vartheta\begin{pmatrix} \alpha \\ \gamma \end{pmatrix}u = \Sigma\, \exp.\left\{J(a\cbtw m)^2 + \tfrac{1}{2}\pi i(m + \alpha)(u + \gamma)\right\},
\]
where the summation extends to all positive and negative even integer values (zero included) of $m$. Here $\alpha$, $\gamma$ are each of them $= 0$ or $1$; and we have the four functions
\[
\vartheta\begin{pmatrix} 0 \\ 0 \end{pmatrix}u,\quad
\vartheta\begin{pmatrix} 1 \\ 0 \end{pmatrix}u,\quad
\vartheta\begin{pmatrix} 0 \\ 1 \end{pmatrix}u,\quad
\vartheta\begin{pmatrix} 1 \\ 1 \end{pmatrix}u,
\]
or, as I prefer to write them,
\[
Au,\quad Bu,\quad Cu,\quad Du.
\]

\noindent\textbf{30.}\quad
The product-theorem gives at once the $16$ equations for the single functions; viz.\ these are the equations expressing the products $\vartheta\begin{pmatrix} \alpha \\ \gamma \end{pmatrix}(u + u') \cdot \vartheta\begin{pmatrix} \alpha' \\ \gamma' \end{pmatrix}(u - u')$ in terms of the functions of $2u$ and $2u'$. The formulae are most conveniently obtained by first expressing the squares and products of the functions of $u$ in terms of the functions of $2u$.

\noindent\textbf{31.}\quad
The formulae for the squares are
\begin{align*}
A^2 u &= A(2u) + B(2u), \\
B^2 u &= A(2u) - B(2u), \\
C^2 u &= C(2u) + D(2u), \\
D^2 u &= C(2u) - D(2u).
\end{align*}

\noindent\textbf{32.}\quad
We then have the product-formulae
\begin{align*}
2Au \cdot Bu &= A'u + B'u, \\
2Cu \cdot Du &= C'u + D'u,
\end{align*}
where $A'$, $B'$, $C'$, $D'$ denote the like functions with a different parameter.

\vspace{1em}

\begin{center}
\textit{Notation by means of radical-functions.}
\end{center}

\vspace{0.5em}

\noindent\textbf{33.}\quad
The foregoing formulae suggest a notation by means of radical-functions. Writing
\begin{align*}
a - x &= \alpha, & b - x &= \beta, & c - x &= \gamma, & d - x &= \delta,
\end{align*}
the radical-functions are
\[
\sqrt{\alpha},\quad \sqrt{\beta},\quad \sqrt{\gamma},\quad \sqrt{\delta},
\]
and the products which correspond to the four functions are
\[
\sqrt{\alpha\beta},\quad \sqrt{\alpha\gamma},\quad \sqrt{\alpha\delta},\quad \sqrt{\beta\gamma},\quad \sqrt{\beta\delta},\quad \sqrt{\gamma\delta}.
\]

\noindent\textbf{34.}\quad
The formulae of the product-theorem then assume a very elegant form; in particular, the equations for the squares become
\begin{align*}
A^2 u &= \mathfrak{A}(a - x), \\
B^2 u &= \mathfrak{A}(b - x), \\
C^2 u &= \mathfrak{C}(c - x), \\
D^2 u &= \mathfrak{D}(d - x),
\end{align*}
where $\mathfrak{A}$, $\mathfrak{B}$, $\mathfrak{C}$, $\mathfrak{D}$ are the radical-functions.

\vspace{1em}

\begin{center}
\textit{Third Part.---The Double Theta-Functions.}
\end{center}

\vspace{0.5em}

\begin{center}
\textit{Definition and fundamental properties.}
\end{center}

\vspace{0.5em}

\noindent\textbf{35.}\quad
As an instance, taking $\begin{matrix} \alpha & \alpha' \\ \gamma & \gamma' \end{matrix} = \begin{matrix} 1 & 0 \\ 1 & 1 \end{matrix}$, the product-equation is
\begin{align*}
\vartheta\begin{pmatrix} 1 \\ 1 \end{pmatrix}(u + u') \cdot \vartheta\begin{pmatrix} 0 \\ 1 \end{pmatrix}(u - u') &= \Theta\begin{pmatrix} \tfrac{1}{2} \\ 0 \end{pmatrix}(2u) \cdot \Theta\begin{pmatrix} \tfrac{1}{2} \\ 0 \end{pmatrix}(2u') + \Theta\begin{pmatrix} \tfrac{1}{2} \\ 1 \end{pmatrix}(2u) \cdot \Theta\begin{pmatrix} \tfrac{1}{2} \\ 1 \end{pmatrix}(2u'), \\
&= iP \cdot P' \qquad\qquad - iQ \cdot Q',
\end{align*}
which agrees with the before-given value.

\noindent\textbf{36.}\quad
The following values are not actually required: but I give them to fix the ideas and to show the meaning of the quantities with which we work.

\begin{center}
\small
\begin{tabular}{l|l}
\multicolumn{1}{c|}{$u = 0$} & \\
\hline
$X = \Theta\begin{pmatrix} 0 \\ 0 \end{pmatrix}(2u) = 1 + 2q^{\frac{1}{2}}\cos 2\pi u + 2q^2\cos 4\pi u + \dots$, & $\alpha = 1 + 2q^{\frac{1}{2}} + 2q^2 + \dots$, \\
$Y = \Theta\begin{pmatrix} 1 \\ 0 \end{pmatrix}(2u) = \phantom{1 +}\ 2q^{\frac{1}{4}}\cos\pi u + 2q^{\frac{9}{4}}\cos 3\pi u + \dots$, & $\beta = \phantom{1 +}\ 2q^{\frac{1}{4}} + 2q^{\frac{9}{4}} + \dots$, \\
$X_1 = \Theta\begin{pmatrix} 0 \\ 1 \end{pmatrix}(2u) = 1 - 2q^{\frac{1}{2}}\cos 2\pi u + 2q^2\cos 4\pi u - \dots$, & $\alpha_1 = 1 - 2q^{\frac{1}{2}} + 2q^2 - \dots$, \\
$Y_1 = \Theta\begin{pmatrix} 1 \\ 1 \end{pmatrix}(2u) = \phantom{1 +}\ 2q^{\frac{1}{4}}\sin\pi u + 2q^{\frac{9}{4}}\sin 3\pi u - \dots$, & $\beta_1' = 2\pi(-q^{\frac{1}{4}} + 3q^{\frac{9}{4}} - \dots)$ \\
 & $\phantom{\beta_1'} = \dfrac{d}{du}Y_1$ for $u = 0$.
\end{tabular}
\end{center}

\begin{align*}
P &= \Theta\begin{pmatrix} \tfrac{1}{2} \\ 0 \end{pmatrix}(2u) = q^{\frac{1}{8}}(\cos\tfrac{1}{2}\pi u + i\sin\tfrac{1}{2}\pi u) + q^{\frac{9}{8}}(\cos\tfrac{3}{2}\pi u - i\sin\tfrac{3}{2}\pi u) \\
  &\qquad\qquad\qquad\qquad + q^{\frac{25}{8}}(\cos\tfrac{5}{2}\pi u + i\sin\tfrac{5}{2}\pi u) + \dots, \\
Q &= \Theta\begin{pmatrix} \tfrac{1}{2} \\ 0 \end{pmatrix}(2u) = q^{\frac{1}{8}}(\cos\tfrac{1}{2}\pi u - i\sin\tfrac{1}{2}\pi u) + q^{\frac{9}{8}}(\cos\tfrac{3}{2}\pi u + i\sin\tfrac{3}{2}\pi u) \\
  &\qquad\qquad\qquad\qquad + q^{\frac{25}{8}}(\cos\tfrac{5}{2}\pi u - i\sin\tfrac{5}{2}\pi u) + \dots, \\[0.5em]
P_1 &= \Theta\begin{pmatrix} \tfrac{1}{2} \\ 1 \end{pmatrix}(2u) = \frac{1 + i}{\sqrt{2}}\{q^{\frac{1}{8}}(\cos\tfrac{1}{2}\pi u + i\sin\tfrac{1}{2}\pi u) - q^{\frac{9}{8}}(\cos\tfrac{3}{2}\pi u - i\sin\tfrac{3}{2}\pi u) \\
  &\qquad\qquad\qquad\qquad - q^{\frac{25}{8}}(\cos\tfrac{5}{2}\pi u + i\sin\tfrac{5}{2}\pi u) + \dots\}, \\[0.5em]
Q_1 &= \Theta\begin{pmatrix} \tfrac{1}{2} \\ 1 \end{pmatrix}(2u) = \frac{1 - i}{\sqrt{2}}\{q^{\frac{1}{8}}(\cos\tfrac{1}{2}\pi u - i\sin\tfrac{1}{2}\pi u) - q^{\frac{9}{8}}(\cos\tfrac{3}{2}\pi u + i\sin\tfrac{3}{2}\pi u) \\
  &\qquad\qquad\qquad\qquad - q^{\frac{25}{8}}(\cos\tfrac{5}{2}\pi u - i\sin\tfrac{5}{2}\pi u) + \dots\};
\end{align*}
and therefore also
\begin{align*}
p = q &= q^{\frac{1}{8}} + q^{\frac{9}{8}} + q^{\frac{25}{8}} + \dots\dots, \\
p_1 = \frac{1 + i}{\sqrt{2}}\{q^{\frac{1}{8}} - q^{\frac{9}{8}} - q^{\frac{25}{8}} + q^{\frac{49}{8}} + q^{\frac{81}{8}} - \dots\}, &\quad
q_1 = \frac{1 - i}{\sqrt{2}}\{\text{Do.}\};\quad p_1 = iq_1.
\end{align*}

\vspace{1em}

\begin{center}
\textbf{480}
\end{center}

\vspace{0.5em}

\noindent\textbf{[704}

\vspace{0.5em}

\begin{center}
\textit{The square set, and $x$-formulae.}
\end{center}

\vspace{0.5em}

\noindent\textbf{37.}\quad
We use the square-set, in the first instance by writing therein $u' = 0$; the equations become
\begin{align*}
A^2 u &= \alpha X + \beta Y, = \omega^2 \mathfrak{A}(a - x), \\
B^2 u &= \beta X + \alpha Y, = \omega^2 \mathfrak{A}(b - x), \\
C^2 u &= \alpha X - \beta Y, = \omega^2 \mathfrak{C}(c - x), \\
D^2 u &= \beta X - \alpha Y, = \omega^2 \mathfrak{D}(d - x);
\end{align*}
viz.\ the equations without their last members show that the squared functions are linear functions of $X$ and $Y$, such that there exist values of $\alpha$, $\beta$, $\omega$; $X$, $Y$ may be assumed equal to the squared functions $\mathfrak{A}(a - x)$, etc., $\omega^2$ being then proportional to $x$; the functions $\mathfrak{A}$, $\mathfrak{B}$, $\mathfrak{C}$, $\mathfrak{D}$ then are proportional to $\mathfrak{A}(a - x)$, etc., so that the squared functions are proportional to $\mathfrak{A}(a - x)$, etc.

To show the meaning of the factor $\omega^2$, observe that, from any two of the equations, for instance from the first and second, we have an equation without $\omega$,
\[
A^2 u - B^2 u = \mathfrak{A}(a - x) - \mathfrak{A}(b - x);
\]
and using this to determine $x$, and then substituting in $\omega^2 = A^2 u \div \mathfrak{A}(a - x)$, we find
\[
\omega^2 = \frac{A^2 u - B^2 u}{(a - b)\mathfrak{A}d},
\]
where the numerator is a function not in anywise more important than any other linear function of $A^2 u$ and $B^2 u$.

\vspace{1em}

\begin{center}
\textbf{481}
\end{center}

\vspace{0.5em}

\noindent\textbf{[704}

\vspace{0.5em}

\noindent\textbf{38.}\quad
The function $Du$ vanishes for $u = 0$, and we may assume that the corresponding value of $x$ is $= d$. Writing in the other equations $u = 0$, they become
\begin{align*}
A^2(0) &= (\alpha^2 + \beta^2) = \omega_0^2 \mathfrak{A}(a - d), \\
B^2(0) &= 2\alpha\beta = \omega_0^2 \mathfrak{A}(b - d), \\
C^2(0) &= (\alpha^2 - \beta^2) = \omega_0^2 \mathfrak{C}(c - d),
\end{align*}
where $\omega_0^2$ is the value of $\omega^2$ for $u = 0$; these equations denote the quantities $\mathfrak{A}$, $\mathfrak{B}$, $\mathfrak{C}$, $\mathfrak{D}$ containing a constant $\omega_0^2$, not absolute equalities but only equalities of ratios: thus, without the $\omega$ factors the last-mentioned equations would denote
\[
A^2(0) : B^2(0) : C^2(0) = \mathfrak{A}(a - d) : \mathfrak{A}(b - d) : \mathfrak{C}(c - d).
\]

The function $Du$ becomes on writing therein $x = d$.

It is convenient to omit the $\omega$ factors and to write; it being understood that, without them, the formulae are
\[
A^2(0) : B^2(0) : C^2(0) = \mathfrak{A}(a - d) : \mathfrak{A}(b - d) : \mathfrak{C}(c - d) : \mathfrak{D}(d - d);
\]
that is,
\[
A^2(0) = \alpha^2 + \beta^2 : 2\alpha\beta : \alpha^2 - \beta^2 = \mathfrak{A}(a - d) : \mathfrak{A}(b - d) : \mathfrak{C}(c - d) : \mathfrak{D}(d - d).
\]

The functions $\mathfrak{A}$, $\mathfrak{B}$, $\mathfrak{C}$, $\mathfrak{D}$ only present themselves in the products $\omega \omega'$, etc., and their absolute magnitudes are therefore essentially indeterminate but regarding $\omega$ as a constant factor of properly determined value, the absolute values of $\mathfrak{A}$, $\mathfrak{B}$, $\mathfrak{C}$, $\mathfrak{D}$ may be regarded as determinate, and this is accordingly done in the sequel.

\vspace{1em}

\begin{center}
\textbf{482}
\end{center}

\vspace{0.5em}

\noindent\textbf{[704}

\vspace{0.5em}

\begin{center}
\textit{Relations between the constants.}
\end{center}

\vspace{0.5em}

\noindent\textbf{39.}\quad
The formulae contain the differences of the quantities $a$, $b$, $c$, $d$; denoting these differences in the usual manner by $f$, $g$, $h$, $a$, $b$, $c$, so that
\begin{align*}
-h + g - a &= 0, & -g + f - b &= 0, & -f + h - c &= 0, \\
a + b + c &= 0,
\end{align*}
and also
\[
af + bg + ch = 0,
\]
and then assuming the absolute value of one of the quantities $\mathfrak{D}$, we have the system of relations
\begin{align*}
\mathfrak{A}^2 &= -agh, & \mathfrak{B}^2 &= bhf, & \mathfrak{C}^2 &= -cfg, & \mathfrak{D}^2 &= abc, \\
\mathfrak{B}\mathfrak{C}a &= -\mathfrak{B}\mathfrak{D}g, & \mathfrak{C}\mathfrak{A}b &= -\mathfrak{C}\mathfrak{D}h, & \mathfrak{A}\mathfrak{B}c &= -\mathfrak{A}\mathfrak{D}f, \\
\mathfrak{B}\mathfrak{C}\mathfrak{D}f &= -\mathfrak{B}gh, & \mathfrak{C}\mathfrak{A}\mathfrak{D}g &= -\mathfrak{C}hf, & \mathfrak{A}\mathfrak{B}\mathfrak{D}h &= -\mathfrak{A}fg, \\
\mathfrak{A}\mathfrak{B}\mathfrak{C}\mathfrak{D} &= abefgh,
\end{align*}
and
\begin{align*}
-a^2\mathfrak{A}^2 + b^2\mathfrak{B}^2 + c^2\mathfrak{C}^2 + f^2\mathfrak{D}^2 &= bcf(af + bg + ch) = 0, \\
a^2\mathfrak{A}^2 - b^2\mathfrak{B}^2 + c^2\mathfrak{C}^2 + g^2\mathfrak{D}^2 &= cag(af + bg + ch) = 0, \\
a^2\mathfrak{A}^2 + b^2\mathfrak{B}^2 - c^2\mathfrak{C}^2 + h^2\mathfrak{D}^2 &= abh(af + bg + ch) = 0, \\
-f^2\mathfrak{A}^2 + g^2\mathfrak{B}^2 + h^2\mathfrak{C}^2 + \mathfrak{D}^2 &= fgh(af + bg + ch) = 0.
\end{align*}

It is to be remarked that, taking $c$, $a$, $b$, $d$ in the order of decreasing magnitude, we have $-a$, $b$, $c$, $f$, $g$, $h$ all positive; hence $\mathfrak{A}^2$, $\mathfrak{B}^2$, $\mathfrak{C}^2$, $\mathfrak{D}^2$ all real; and taking as we may do, $\mathfrak{D}$ negative, then $\mathfrak{A}$, $\mathfrak{B}$, $\mathfrak{C}$ may be taken positive; that is, we have $-a$, $b$, $c$, $g$, $h$, $\mathfrak{A}$, $\mathfrak{B}$, $\mathfrak{C}$, $-\mathfrak{D}$ all of them positive.

\noindent\textbf{40.}\quad
We have
\begin{align*}
A^2(0) &= \alpha^2 + \beta^2 = \mathfrak{A}f, \\
B^2(0) &= 2\alpha\beta = \mathfrak{B}g, \\
C^2(0) &= \alpha^2 - \beta^2 = \mathfrak{C}h.
\end{align*}

The foregoing equations

\vspace{1em}

\begin{center}
\textbf{483}
\end{center}

\vspace{0.5em}

\noindent\textbf{[704}

\vspace{0.5em}

\noindent\textbf{41.}\quad
We have $A^2 u$, $B^2 u$, $C^2 u$, $D^2 u$ proportional to $\mathfrak{A}(a - x)$, $\mathfrak{A}(b - x)$, $\mathfrak{C}(c - x)$, $\mathfrak{D}(d - x)$ respectively; or, what is the same thing, $\sqrt{A}$, $\sqrt{B}$, $\sqrt{C}$, $\sqrt{D}$ are proportional to the radical-functions. And hence conversely, the variable $x$ is determined as a rational function of $A^2 u$, $B^2 u$, $C^2 u$, $D^2 u$; viz.
\[
x = \frac{a\mathfrak{A}^2 A^2 u + b\mathfrak{B}^2 B^2 u + c\mathfrak{C}^2 C^2 u + d\mathfrak{D}^2 D^2 u}{\mathfrak{A}^2 A^2 u + \mathfrak{B}^2 B^2 u + \mathfrak{C}^2 C^2 u + \mathfrak{D}^2 D^2 u}.
\]

\noindent\textbf{42.}\quad
The differential relation between $x$ and $u$ is
\[
\frac{dx}{\sqrt{a - x . b - x . c - x . d - x}} = M\, du,
\]
where $M$ is a constant; and the integral of this equation, determining $u$ as a function of $x$, is the inverse function which gives rise to the whole theory.

\noindent\textbf{43.}\quad
The foregoing formulae establish the connexion between the theta-functions and the elliptic functions; and they serve to express the latter in terms of the former. In particular, we have
\[
\sn Ku = \sqrt{\frac{\mathfrak{D}}{\mathfrak{C}}} \cdot \frac{Du}{Cu},\quad
\cn Ku = \sqrt{\frac{\mathfrak{D}}{\mathfrak{B}}} \cdot \frac{Bu}{Cu},\quad
\dn Ku = \sqrt{\frac{\mathfrak{C}}{\mathfrak{A}}} \cdot \frac{Au}{Cu}.
\]

\vspace{1em}

\begin{center}
\textit{The product-equations.}
\end{center}

\vspace{0.5em}

\noindent\textbf{44.}\quad
The $16$ product-equations for the double theta-functions may be written in the form:

For the square-set:
\begin{align*}
A(u + u') \cdot A(u - u') &= AA' + BB' + CC' + DD', \\
B(u + u') \cdot B(u - u') &= AB' + BA' + CD' + DC', \\
C(u + u') \cdot C(u - u') &= AC' + BD' + CA' + DB', \\
D(u + u') \cdot D(u - u') &= AD' + BC' + CB' + DA'.
\end{align*}

\noindent\textbf{45.}\quad
And for the other products:
\begin{align*}
AB(u + u') \cdot AB(u - u') &= \text{etc.}
\end{align*}

\vspace{1em}

\begin{center}
\textit{The addition-theorem.}
\end{center}

\vspace{0.5em}

\noindent\textbf{46.}\quad
The addition-theorem is obtained from the product-equations by eliminating the functions of $(u - u')$ between any two of the equations. The resulting formula expresses the quotient-functions of $(u + u')$ in terms of the quotient-functions of $u$ and $u'$ respectively.

\vspace{1em}

\begin{center}
\textbf{484--491}
\end{center}

\vspace{0.5em}

\noindent\textbf{[704}

\vspace{0.5em}

\begin{center}
\textit{Addition-formulae.}
\end{center}

\vspace{0.5em}

\noindent\textbf{56.}\quad
The addition-theorem of the quotient-functions for the elliptic functions is of course obtained in my paper ``On the Double $\vartheta$-Functions'' (Crelle, t.~LXXXVII.\ (1879), pp.~74---81, [697]); viz.
\[
dx = \frac{dy}{\sqrt{Y}} = \frac{dz}{\sqrt{Z}},
\]
for the differential equation
\[
\frac{dx}{\sqrt{X}} = \frac{dy}{\sqrt{Y}} = \frac{dz}{\sqrt{Z}},
\]
the particular integral of which to obtain in the formulae of the paper just referred to, reduces itself to $z = x$, and writing for shortness $a_1 = a - y$, $b_1 = b - y$, $c_1 = c - y$, $d_1 = d - y$, then when the interchange is made, the formulae become

\begin{align*}
\sqrt{d - b . d - c}\, \{&\sqrt{adb_1 c_1} + \sqrt{a_1 d_1 bc}\} \\
&= (bc, ad)\, \sqrt{d - b . d - c}\, \frac{X - Y}{\sqrt{adb_1 c_1} - \sqrt{a_1 d_1 bc}}, \\
\sqrt{d - b . d - c}\, \{&\sqrt{bdc_1 a_1} + \sqrt{b_1 d_1 ca}\} \\
&= (ac, bd)\, \sqrt{d - b . d - c}\, \frac{X - Y}{\sqrt{bdc_1 a_1} - \sqrt{b_1 d_1 ca}}, \\
\sqrt{d - b . d - c}\, \{&\sqrt{cda_1 b_1} + \sqrt{c_1 d_1 ab}\} \\
&= (ab, cd)\, \sqrt{d - b . d - c}\, \frac{X - Y}{\sqrt{cda_1 b_1} - \sqrt{c_1 d_1 ab}},
\end{align*}
where $(bc, ad)$, $(ac, bd)$, $(ab, cd)$ denote respectively
\[
\begin{vmatrix}
1, & x + y, & xy \\
1, & b + c, & bc \\
1, & a + d, & ad
\end{vmatrix},\quad
\begin{vmatrix}
1, & x + y, & xy \\
1, & a + c, & ac \\
1, & b + d, & bd
\end{vmatrix},\quad
\begin{vmatrix}
1, & x + y, & xy \\
1, & a + b, & ab \\
1, & c + d, & cd
\end{vmatrix}.
\]

\vspace{1em}

\begin{center}
\textbf{492}
\end{center}

\vspace{0.5em}

\noindent\textbf{[704}

\vspace{0.5em}

\noindent\textbf{57.}\quad
In the foregoing formulae, $(bc, ad)$, $(ac, bd)$, and $(ab, cd)$ denote respectively
\[
\begin{vmatrix}
1, & x + y, & xy \\
1, & b + c, & bc \\
1, & a + d, & ad
\end{vmatrix},\quad
\begin{vmatrix}
1, & x + y, & xy \\
1, & a + c, & ac \\
1, & b + d, & bd
\end{vmatrix},\quad
\begin{vmatrix}
1, & x + y, & xy \\
1, & a + b, & ab \\
1, & c + d, & cd
\end{vmatrix};
\]
and substituting for $\mathfrak{A}$, $\mathfrak{B}$, $\mathfrak{C}$, $\mathfrak{D}$, $b - x$, etc., we have moreover
\begin{align*}
A^2(u + v) &= \sqrt{c - b . c - d . c - a}\, (a - x)(a - y)(a - z), \\
B^2(u + v) &= \sqrt{c - a . c - d . a - d}\, (b - x)(b - y)(b - z), \\
C^2(u + v) &= \sqrt{a - b . a - d . b - d}\, (c - x)(c - y)(c - z), \\
D^2(u + v) &= \sqrt{c - b . c - a . a - b}\, (d - x)(d - y)(d - z),
\end{align*}
the constant multipliers being of course the same in the three columns.

According to what precedes, the radical of the fourth line should be taken with the sign $-$.

The functions $(bc, ad)$, etc., contained in the formulae, require a transformation such as
\[
(b - c)(bc, ad) = \begin{vmatrix} b - x, & b - y \\ c - x, & c - y \end{vmatrix} \div \begin{vmatrix} b - a, & b - d \\ c - a, & c - d \end{vmatrix}
\]
in order to make them separately homogeneous in the differences $a - x$, $b - x$, $c - x$, $d - x$, and $a - y$, $b - y$, $c - y$, $d - y$, and therefore to make them expressible as linear functions of the squared functions $Au$, etc., and $A'u$, etc., respectively; the formulae then give the quotient-functions $\dfrac{A(u + v)}{D(u + v)}$, etc., in terms of the quotient-functions of $u$ and $v$ respectively.

\vspace{1em}

\begin{center}
\textit{Doubly infinite product-forms.}
\end{center}

\vspace{0.5em}

\noindent\textbf{58.}\quad
The product-functions $Au$, $Bu$, $Cu$, $Du$ may be expressed each as a doubly infinite product. Writing for shortness
\begin{align*}
m + \tfrac{1}{2} &= (m, n), & m - \tfrac{1}{2} &= (m, n), \\
n + \tfrac{1}{2} &= (m, n), & n - \tfrac{1}{2} &= (m, n),
\end{align*}
the formulae are
\begin{align*}
Au &= \prod_{m,n} \left(1 - \frac{u}{(m,n)}\right)\left(1 - \frac{u}{(m,n)}\right)\left(1 - \frac{u}{(m,n)}\right)\left(1 - \frac{u}{(m,n)}\right), \\
Bu &= \text{etc.},
\end{align*}
with the proper constant factors.

\vspace{1em}

\begin{center}
\textbf{493--500}
\end{center}

\vspace{0.5em}

\noindent\textbf{[704}

\vspace{0.5em}

\noindent\textbf{59.}\quad
The several formulae lead to the expression of the functions in terms of the parameters $a$, $b$, $c$, $d$, and the arguments $u$, $v$; and they establish the connexion with the hyperelliptic functions.

\noindent\textbf{60.}\quad
The foregoing theory is carried out in detail in the memoir; and it is shown how the functions may be calculated numerically for given values of the parameters and arguments. The numerical calculations are illustrated by examples; and the results are compared with those obtained by other methods.

\noindent\textbf{61.}\quad
The memoir concludes with some remarks on the extension of the theory to the general case of $p$ arguments; and on the application of the theta-functions to the solution of problems in the theory of Abelian functions.

\vspace{1em}

\begin{center}
\rule{0.5\textwidth}{0.4pt}
\end{center}

\vspace{1em}

\begin{center}
\textit{End of Papers 702, 703, 704}
\end{center}

\vspace{1em}

\begin{flushright}
\textit{C. X.}
\end{flushright}

\end{document}
